\section{Методология измерения}
\label{sec:method}

\subsection{Конфигурации исполнения}
\label{sec:sources}

Для ответа на вопрос об источнике расхождений построены четыре
конфигурации исполнения одной модели (рис.~\ref{fig:configs}).
В CPU-эталоне модель исполняется средствами PyTorch без XLA.
В трёх остальных конфигурациях используется симулятор и общая цепочка
компиляции OpenXLA~\cite{OpenXLA2023}: модель преобразуется в граф
высокоуровневых операций (High-Level Operations, HLO), после чего
выбирается одна из реализаций GEMM. Это последовательное суммирование
на CPU, cuBLAS~\cite{cuBLAS2024} или CUDA-ядро с явно заданным
блочным порядком накопления~\cite{CUDA2024}.

Сравнение (a) связывает CPU-эталон с симулятором, в котором GEMM
исполняется на CPU. Оно охватывает изменение всего пути исполнения.
В сравнениях (b) и (d) CPU-реализация GEMM заменяется соответственно
на cuBLAS и CUDA-ядро; сравнение (c) сопоставляет эти GPU-реализации
между собой. Таким образом, в парах (b)--(d) при одинаковых входах слоя сохраняется
общая цепочка компиляции и проверяется замена GEMM.

\begin{figure}[htbp]
    \centering
    \begin{tikzpicture}[
        font=\normalsize,
        box/.style={draw=black!65, rounded corners=2pt, fill=black!4,
            text width=25mm, minimum height=9mm, align=center},
        arr/.style={-{Stealth[length=4pt]}, line width=0.6pt}
    ]
        \node[box] (model) at (0,0) {Модель\\PyTorch};
        \node[box, fill=babyblue!20] (ref) at (0,-1.5) {CPU-эталон\\без XLA};
        \node[box] (xla) at (3.4,0) {XLA\\граф HLO};
        \node[box] (sim) at (6.6,0) {Симулятор\\GEMM};
        \node[box, fill=babyblue!20] (cpu) at (10,1.35) {CPU};
        \node[box, fill=green!8] (blas) at (10,0) {cuBLAS};
        \node[box, fill=orange!12] (cuda) at (10,-1.35) {CUDA-ядро};
        \draw[arr] (model) -- (ref);
        \draw[arr] (model) -- (xla);
        \draw[arr] (xla) -- (sim);
        \coordinate (fork) at (8.35,0);
        \draw[line width=0.6pt] (sim.east) -- (fork);
        \draw[line width=0.6pt] (8.35,-1.35) -- (8.35,1.35);
        \draw[arr] (8.35,1.35) -- (cpu.west);
        \draw[arr] (fork) -- (blas.west);
        \draw[arr] (8.35,-1.35) -- (cuda.west);
        \fill (fork) circle (1pt);
        \node[anchor=north, align=left] at (4.7,-2.2) {
            (a) эталон / XLA+CPU\qquad (b) XLA+CPU / cuBLAS\\
            (c) cuBLAS / CUDA\qquad (d) XLA+CPU / CUDA
        };
    \end{tikzpicture}
    \caption{Пути исполнения одной модели.
        Стрелки показывают передачу вычислений; под схемой указаны
        пары сравнений. Для пар (b)--(d) цепочка компиляции OpenXLA общая.}
    \label{fig:configs}
\end{figure}

Общая цепочка контролирует компиляцию в парах (b)--(d),
но разность норм расхождений не выделяет её вклад:
последующие нелинейные операции преобразуют возникшие ошибки.
Отдельные механизмы компиляции здесь не измеряются;
формат входов \texttt{FP32} не задаёт формат промежуточных операций.

\subsection{Изолированное измерение и распространение}

Даже при фиксированной паре реализаций сравнение полных запусков
не показывает, в каком слое возникло отклонение: изменённый вход может
влиять на все последующие результаты. Чтобы выделить собственную
ошибку слоя, сначала сохраняются тензоры опорного запуска, затем каждый
слой сравниваемой реализации повторяется на них независимо.
В обратном проходе фиксируются как тензоры
прямого прохода, так и опорный градиент по выходу узла, включающий
сумму вкладов всех его потребителей. Ошибки других сравниваемых узлов
в такой повтор не передаются.

Тензоры сохраняются средствами перехвата PyTorch~\cite{Paszke2019}.
При отложенном исполнении XLA первый проход строит граф и фиксирует
точки захвата, а второй извлекает вычисленные значения для сравнения.



После изолированного измерения проверяется совместное действие
сохранённых локальных разностей. Обозначим их через
$\widehat\delta_l$ для прямого и $\widehat\eta_x^{(j)}$ для обратного
прохода. Шляпка отмечает значения, измеренные на опорных входах.
В прямом испытании разность прибавляется к выходу опорного узла перед
передачей всем его потребителям. В обратном испытании она прибавляется
к вычисленному градиенту по активациям, а тензоры прямого прохода
остаются опорными. Схема на рис.~\ref{fig:injection} показывает
отдельный шаг каждого из двух раздельных запусков.
Полные матрицы якобианов при таком испытании строить не требуется.

\begingroup
\setlength{\intextsep}{3pt}
\begin{figure}[H]
    \centering
    \begin{tikzpicture}[
        font=\normalsize,
        op/.style={draw=black!65, rounded corners=2pt,
            fill=babyblue!20, text width=25mm, minimum height=9mm, align=center},
        add/.style={draw=black!65, circle, minimum size=6mm, inner sep=1pt},
        arr/.style={-{Stealth[length=4pt]}, line width=0.6pt},
        loc/.style={arr, red!70!black}
    ]
        \node[align=left] at (-1.4,0.9) {Прямой\\запуск};
        \node (x) at (0,0) {$\widetilde{\Vec{x}}^{(l)}$};
        \node[op] (forward) at (2.6,0) {Опорный\\узел $l$};
        \node[add] (sumf) at (5.6,0) {$+$};
        \node (y) at (8.3,0) {$\widetilde{\Vec{y}}^{(l)}$};
        \node (delta) at (5.6,1.1) {$\widehat\delta_l$};
        \draw[arr] (x) -- (forward);
        \draw[arr] (forward) -- (sumf);
        \draw[arr] (sumf) -- (y);
        \draw[loc] (delta) -- (sumf);
        \node[align=left] at (-1.4,-1.4) {Обратный\\запуск};
        \node[align=center] (gin) at (8.3,-2.3)
            {Вклады\\$\widetilde{\Vec{g}}^{(k)}$\\$k\in\mathcal{S}_j$};
        \node[op] (backward) at (5.6,-2.3) {Опорный\\шаг};
        \node[add] (sumb) at (2.6,-2.3) {$+$};
        \node (gout) at (0,-2.3) {$\widetilde{\Vec{g}}^{(j)}$};
        \node (eta) at (2.6,-1.25) {$\widehat\eta_x^{(j)}$};
        \draw[arr] (gin) -- (backward);
        \draw[arr] (backward) -- (sumb);
        \draw[arr] (sumb) -- (gout);
        \draw[loc] (eta) -- (sumb);
    \end{tikzpicture}
    \caption{Введение сохранённых локальных разностей после опорных
        вычислений. Ряды соответствуют отдельным испытаниям.
        В обратном запуске тензоры прямого прохода фиксированы.}
    \label{fig:injection}
\end{figure}
\endgroup

Для прямого прохода отношение итогового и максимального локального
расхождений задаётся~\eqref{eq:ratio}.
В показанном обратном режиме используются опорные тензоры прямого
прохода~\eqref{eq:bwd}; обозначение $R_{\mathrm{bwd}}$ относится к
аналогичному отношению расхождений градиентного выхода и локальных градиентов.
Дополнительный вклад изменённого прямого прохода разделяется
по~\eqref{eq:grad-decomp}. Измерения при фиксированном входе
не устанавливают свойства ансамбля.

\subsection{Компоненты и условия эксперимента}

Чтобы проверить зависимость переноса от структуры сети, используются
не только отдельные GEMM, но и исполняемые компоненты нейросетевых
архитектур. Под компонентом здесь понимается сеть или блок с заданной
схемой механизма внимания, нормализации и полносвязных преобразований.
Набор из 14 таких конфигураций приведён в табл.~\ref{tab:components}.
Позиционные варианты обозначаются как rotary position embedding (RoPE)
и attention with linear biases (ALiBi).
В блоках SwiGLU и GeGLU две линейные ветви перемножаются поэлементно,
причём одна предварительно преобразуется нелинейной функцией.
В описании механизма внимания $Q$, $K$, $V$~--- матрицы запросов, ключей
и значений. Скалярные произведения запросов и ключей задают оценки
их сходства, по которым формируются веса для объединения значений.
Полносвязный блок далее обозначается как feed-forward network (FFN).
LayerNorm сокращённо обозначается LN; Pre-Norm~--- нормализацию перед
преобразованием подслоя, Post-LN~--- после остаточного сложения.

Размерности трансформерных компонентов выбраны по пропорциям
Llama~3~8B~\cite{Dubey2024} и Mistral~7B~\cite{Jiang2023}.
Ширина уменьшена в четыре раза: скрытая размерность $4096\to1024$,
размерность FFN $14336\to3584$, число голов $32\to8$;
размерность головы 128 сохранена. Тем самым отношение размерности FFN
к скрытой размерности остаётся равным $3.5$.
Для испытаний выбраны два блока, длина последовательности 32
и размер пакета 2. В варианте со сгруппированными запросами
(grouped-query attention, GQA) число голов ключей и значений уменьшено
с восьми до двух, с сохранением отношения $1/4$ к числу голов запросов.

Входы генерируются из $\mathcal{N}(0,1)$ при фиксированной
инициализации генератора. Начальные веса линейных проекций создаются
реализацией инициализации слоёв PyTorch~\cite{Paszke2019}:
они равномерно распределены на
$[-1/\sqrt{n_{\mathrm{in}}},1/\sqrt{n_{\mathrm{in}}}]$,
где $n_{\mathrm{in}}$~--- число входных активаций.
Масштабы нормализации задаются равными единице.
Отдельно заданные параметры инициализации сохраняются:
для LeNet используется равномерное распределение на $[-2,2]$.

Для обратного прохода используется перекрёстная энтропия
(\texttt{CrossEntropyLoss}) с отдельной выходной линейной проекцией
на 32000 координат. Эта проекция задаёт одинаковое число выходных
классов для всех компонентов; 32000 соответствует размеру словаря
Mistral~7B~\cite{Jiang2023}. Размер испытательной головы фиксируется
отдельно от масштабирования ширины компонентов.
Стандартное отклонение нормальной инициализации её весов равно $0.02$.
Это компонентные испытания на инициализированных параметрах,
а не измерения качества обученных полноглубинных моделей.

Для DeepNorm масштабирование инициализации весов подслоёв из исходной
работы~\cite{Wang2022} не применяется. Поэтому компонент воспроизводит
исследуемую схему связей, но не весь протокол DeepNet.
GPU-реализации исполнялись на NVIDIA A100 SXM4 с 80~ГБ памяти;
CPU-реализация~--- на том же вычислительном узле.
При оценке ускорения учитываются прямой и обратный проходы,
но не компиляция XLA. Отношение рассчитывается по показаниям симулятора,
а не по полному времени работы процесса.
\begin{table}[!htbp]
    \centering
    \caption{Набор из 14 исследуемых компонентов нейросетевых архитектур.}
    \label{tab:components}
    \fitwidth{%
        \begin{tabular}{l l l}
            \hline
            \textbf{Группа}  & \textbf{Компоненты}                & \textbf{Варьируемый фактор}               \\
            \hline
            Базовые (2)      & LeNet                              & свёртки, сигмоида, усреднение        \\
                             & RoBERTa                            & \shortstack[l]{Pre-Norm LayerNorm,\\двунаправленный механизм внимания}    \\
            \hline
            \shortstack[l]{Механизм\\внимания (8)}     & RoPE, ALiBi                        & позиционное кодирование                   \\
                             & GQA                                & число голов K/V                \\
                             & Кеш ключей и значений, скользящее окно            & переиспользование ключей и значений; окно \\
                             & FlashAttention                    & алгоритм механизма внимания              \\
                             & Нормализация Q/K, Softcapping               & нормализация запросов и ключей; ограничение оценок                      \\
            \hline
            FFN (2)          & SwiGLU, GeGLU                      & нелинейная функция одной из ветвей             \\
            \hline
            Схема блока (2) & RMSNorm FFN                        & Pre-Norm RMSNorm (без вычитания среднего) \\
                             & DeepNorm                           & масштабированный Post-LN                  \\
            \hline
        \end{tabular}%
    }
\end{table}
